% Data Archiving System TRS; WBS 123456.02.01.
\section{Technical Requirements Specification Purpose and Description}

\subsection{Purpose}

The purpose of this Technical Requirements Specification (TRS) is to define the technical requirements and expected capabilities of the Data Archiving System (DAS) within the Early Warning Alarm System (EWAS) programme.

The DAS is responsible for the storage, organization, preservation, management and retrieval of mission data generated or received during EWAS operations. It provides the archival infrastructure required to maintain both original mission data and processed data products throughout the operational lifetime of the system and for subsequent historical analysis.

The specification defines the required functions, interfaces, operational characteristics and constraints of the DAS, including its two main elements: the Raw Data Archive (RDA) and the Processed Data Archive (PDA).

The requirements specified in this document are intended to provide a technical basis for the design, implementation, verification and acceptance of the DAS. They are defined independently of specific hardware, software or storage technologies to allow for alternative technical solutions.

\subsection{Scope}

The scope of this specification covers the Data Archiving System as a component of the EWAS Ground Segment. It includes the requirements applicable to the storage and management of raw and processed mission data.

The DAS comprises the following elements:

\begin{itemize}
    \item \textbf{Raw Data Archive (RDA):} responsible for receiving, storing and preserving original, unprocessed data obtained from the Space Segment and Ground Segment.
    \item \textbf{Processed Data Archive (PDA):} responsible for storing and managing processed data products generated through mission data processing, analysis and fusion activities.
\end{itemize}

The scope of the DAS includes:

\begin{itemize}
    \item reception and ingestion of mission data from external system elements;
    \item organization and storage of raw and processed data;
    \item storage of event labels and metadata required for the future training, validation and evaluation of detection and classification algorithms;
    \item management of data identifiers and associated metadata;
    \item preservation of data integrity and traceability;
    \item retrieval and provision of archived data to authorized users and system elements;
    \item management of data access permissions;
    \item backup and recovery of archived data;
    \item support for historical data analysis and future mission-related activities.
\end{itemize}

The DAS is located within the ground-based infrastructure of the EWAS Ground Segment.

The scope of this specification does not include the generation of raw observations by spacecraft, the processing algorithms responsible for producing mission-level analytical results, or the operational decision-making performed by authorized authorities. These activities are considered external functions that provide data to or retrieve data from the DAS.

The specification does not prescribe a particular database management system, storage medium, software platform or physical infrastructure. Such implementation decisions are outside the scope of this document unless explicitly required by an approved system-level constraint.

\subsection{System Overview}

The Data Archiving System provides a centralized archival capability for mission data generated and received during EWAS operations. It supports the preservation and retrieval of information required for operational activities, mission analysis, historical reference and future improvement of detection and assessment processes.

The DAS is divided into two complementary archival elements: the Raw Data Archive and the Processed Data Archive. This separation ensures that original acquired data remain distinguishable from information produced through subsequent processing and analysis.

\subsubsection*{Data Archiving System (DAS)}

The DAS operates as part of the EWAS Ground Segment and interfaces with relevant mission elements responsible for data acquisition, transmission, processing and operational use.

Its primary purpose is to maintain an organized and reliable repository of mission data. The system shall support the ingestion, storage, identification, retrieval and controlled access to archived information.

The DAS shall maintain associations between related raw and processed data products to support traceability from original observations to derived mission information.

The system shall provide access to archived data to authorized user groups, including:

\begin{itemize}
    \item Ground Segment operators responsible for monitoring and operating the mission infrastructure;
    \item authorized national crisis-management and emergency-response authorities requiring access to relevant mission information;
    \item analysts responsible for mission data assessment and historical analysis;
    \item technical and maintenance personnel responsible for maintaining the DAS infrastructure.
\end{itemize}

Access to data shall be controlled according to the user's authorization level and assigned responsibilities.

\subsubsection*{Raw Data Archive (RDA)}

The Raw Data Archive is responsible for preserving original, unprocessed mission data received from the Space Segment and Ground Segment.

The RDA shall maintain the original content of received data without replacing it with processed or interpreted information. It shall support the identification and organization of archived data using relevant metadata, including source, acquisition time, data type and associated event identifiers, where available.

The RDA is intended to store data such as:

\begin{itemize}
    \item telemetry received from spacecraft;
    \item sensor observations and initial acquisition data;
    \item data received from ground stations;
    \item other original mission data required for operational or historical purposes.
\end{itemize}

The RDA shall support data integrity verification and preserve sufficient information to establish the origin and acquisition context of stored data.

Archived raw data shall remain available for authorized retrieval and may serve as input for subsequent processing, verification and historical analysis activities.

\subsubsection*{Processed Data Archive (PDA)}

The Processed Data Archive is responsible for storing and managing mission data products generated through processing, analysis and fusion activities performed by the EWAS system.

The PDA shall maintain processed information together with relevant metadata describing its origin, processing status and association with source data.

The PDA is intended to store data products such as:

\begin{itemize}
    \item candidate event records and assessed event information;
    \item preliminary and refined trajectory estimates;
    \item probabilistic impact area estimates;
    \item preliminary warning information and related status records;
    \item results of mission data analysis and fusion;
    \item historical data products selected for future analysis and reference.
\end{itemize}

The PDA shall support the association of processed data products with corresponding raw data, where such relationships are available.

The archive shall preserve information necessary to identify the processing history of stored products, including relevant timestamps, source references and processing identifiers or versions, where applicable.

The PDA shall provide authorized users and system elements with the ability to retrieve relevant processed information for operational support, historical analysis and mission assessment.

\subsection{Operational Environment}

The DAS operates within the ground-based infrastructure of the EWAS Ground Segment. It receives mission data from relevant system elements and provides archival services to authorized users and connected systems.

The operational environment includes:

\begin{itemize}
    \item the Space Segment, which generates observations and telemetry transmitted to the Ground Segment;
    \item Ground Segment elements responsible for receiving, processing and distributing mission data;
    \item the DAS, which stores and manages raw and processed mission data;
    \item authorized operational users and external authorities requiring access to selected archived information;
    \item technical and maintenance personnel responsible for system administration and maintenance.
\end{itemize}

The DAS shall operate in coordination with the relevant Ground Segment functions and shall support the exchange of data through defined interfaces.

The system shall account for operational constraints associated with data transmission interruptions, temporary loss of connectivity, incomplete data reception and possible failures of storage or supporting infrastructure.

The DAS shall support the preservation of already archived information during temporary communication or service interruptions. Where data cannot be received or stored successfully, the system shall provide appropriate status information to the relevant Ground Segment functions.

\subsection{System Boundaries and Interfaces}

The DAS is a dedicated archival subsystem within the EWAS Ground Segment. Its boundaries are defined by the functions related to data ingestion, storage, management, retrieval and controlled access.

The main interfaces of the DAS include:

\begin{itemize}
    \item \textbf{Space Segment interface:} provides access to mission data originating from spacecraft through the Ground Segment communication and reception infrastructure.
    \item \textbf{Ground Segment interface:} enables the reception of raw mission data and processed data products generated by Ground Segment functions.
    \item \textbf{Operational user interface:} enables authorized operators, analysts and other permitted users to search, retrieve and access archived information.
    \item \textbf{External authority interface:} supports the provision of selected mission data products to authorized national crisis-management and emergency-response authorities, subject to applicable access permissions.
    \item \textbf{Maintenance and administration interface:} supports system administration, monitoring, maintenance, backup and recovery activities.
\end{itemize}

The detailed technical characteristics and data exchange formats of these interfaces shall be defined in the relevant interface documentation.

\subsection{Data Lifecycle and Management}

The DAS shall support the management of mission data throughout the archival lifecycle, from reception to storage, retrieval and eventual disposition in accordance with the applicable data retention policy.

The main data management activities include:

\begin{itemize}
    \item \textbf{Data ingestion:} reception, identification and registration of incoming mission data.
    \item \textbf{Data validation:} verification of data completeness, format and integrity to the extent supported by the available source information.
    \item \textbf{Data classification:} assignment of data to the appropriate archive according to its status as raw or processed mission data.
    \item \textbf{Data storage:} organization and preservation of data within the RDA or PDA.
    \item \textbf{Metadata management:} maintenance of information required to identify, search, associate and retrieve archived data.
    \item \textbf{Data retrieval:} provision of archived information in response to authorized requests.
    \item \textbf{Data backup and recovery:} creation and maintenance of backup copies and restoration of data following failures or data loss.
    \item \textbf{Data retention:} preservation and eventual disposition of data according to an approved retention policy.
\end{itemize}

The minimum retention period for raw data stored in the RDA shall be 10 years. The minimum retention period for processed data stored in the PDA shall also be 10 years.

Data relating to detected events, false alarms, missed events, anomalies and warnings shall be preserved for the entire duration of the EWAS programme. Metadata, processing history and algorithm identifiers and versions shall be retained for at least as long as the data products to which they relate. Longer or indefinite preservation may be applied to mission-critical data where permitted by applicable legal, security and resource constraints.

The DAS shall maintain the distinction between original raw data and derived processed data throughout the archival lifecycle. Any authorized modification, transformation or replacement of archived information shall be traceable.

\subsection{Historical Data Analysis and Algorithm Improvement}

The DAS shall preserve data and associated information required to support historical analysis and the future training, validation and evaluation of detection and classification algorithms. The DAS shall archive such information but shall not perform model training or algorithm improvement, which remain functions of the relevant Data Processing System or other authorized external functions.

For data products used in historical analysis or algorithm improvement, the DAS shall preserve, where applicable:

\begin{itemize}
    \item event labels and Ground Truth information identifying whether an event was confirmed, false, missed or ambiguous;
    \item references linking processed data products and results to the corresponding source data;
    \item identifiers and versions of the algorithms used to generate or analyse the data;
    \item the time at which processing or analysis was performed;
    \item the confidence level and uncertainty associated with the result;
    \item records of anomalies, processing errors, data-quality issues and other relevant irregularities.
\end{itemize}

The preserved information shall support the reconstruction of the relationship between source observations, processing results, event assessments and warnings. This traceability shall enable the analysis of false positives, false negatives, erroneous trajectory estimates, communication or processing failures, and other system-performance issues.

\subsection{Security and Data Protection}

The DAS shall provide mechanisms to protect archived mission data against unauthorized access, modification, disclosure and loss.

Security provisions shall include:

\begin{itemize}
    \item access control based on user authorization and assigned responsibilities;
    \item protection of data during storage and exchange through appropriate security measures;
    \item preservation and verification of data integrity;
    \item maintenance of records of relevant access and administrative activities;
    \item backup and recovery mechanisms supporting data availability;
    \item protection of archival services against unauthorized operations.
\end{itemize}

The DAS shall distinguish between user access required for operational activities, analytical activities and system maintenance.

Access to mission information provided to external authorities shall be limited to the data and functions authorized for the respective recipients.

The detailed security architecture and applicable security measures shall be defined in coordination with the relevant EWAS Ground Segment and programme-level security requirements.

\subsection{System Constraints and Assumptions}

The DAS design and implementation shall take into account the following constraints and assumptions:

\begin{itemize}
    \item the DAS forms part of the EWAS Ground Segment and operates within its ground-based infrastructure;
    \item the system comprises separate archival elements for raw and processed mission data;
    \item data are received from external mission elements through defined Ground Segment interfaces;
    \item the DAS is responsible for archival functions and does not replace the processing or operational decision-making functions of other EWAS elements;
    \item access to archived information is restricted to authorized users and systems;
    \item raw data in the RDA and processed data in the PDA shall each be retained for a minimum of 10 years;
    \item data concerning detected events, false alarms, missed events, anomalies and warnings shall be retained for the entire duration of the EWAS programme;
    \item metadata, processing history and algorithm identifiers and versions shall be retained for at least as long as the data products to which they relate;
    \item the implementation shall remain independent of specific hardware and software technologies unless otherwise constrained by approved programme-level requirements.
\end{itemize}

The DAS shall be designed to support the operational and archival needs of EWAS while maintaining data integrity, traceability, availability and controlled access throughout its defined service lifetime.